\section{The generating function of a weighted set}

\subsection{The theory}

\begin{definition}
\label{def.gf-ws.weighted-sets}
\lean{AlgebraicCombinatorics.WeightedSet, AlgebraicCombinatorics.WeightedSet.IsFiniteType, AlgebraicCombinatorics.WeightedSet.weightGenFun, AlgebraicCombinatorics.WeightedSet.Isomorphism, AlgebraicCombinatorics.WeightedSet.AreIsomorphic}
\leantarget
\leanok
\textbf{(a)} A \emph{weighted set} is a set $A$
equipped with a function $w:A\rightarrow\mathbb{N}$, which is called the
\emph{weight function} of this weighted set. For each $a\in A$, the value
$w\left(  a\right)  $ is denoted $\left\vert a\right\vert $ and is called the
\emph{weight} of $a$ (in our weighted set). \medskip

Usually, instead of explicitly specifying the weight function $w$ as a
function, we will simply specify the weight $\left\vert a\right\vert $ for
each $a\in A$. The weighted set consisting of the set $A$ and the weight
function $w$ will be called $\left(  A,w\right)  $ or simply $A$ when the
weight function $w$ is clear from the context. \medskip

\textbf{(b)} A weighted set $A$ is said to be \emph{finite-type} if for each
$n\in\mathbb{N}$, there are only finitely many $a\in A$ having weight
$\left\vert a\right\vert =n$. \medskip

\textbf{(c)} If $A$ is a finite-type weighted set, then the \emph{weight
generating function} of $A$ is defined to be the FPS%
\[
\sum_{a\in A}x^{\left\vert a\right\vert }=\sum_{n\in\mathbb{N}}\left(
\text{\# of }a\in A\text{ having weight }n\right)  \cdot x^{n}\in
\mathbb{Z}\left[  \left[  x\right]  \right]  .
\]
This FPS is denoted by $\overline{A}$. \medskip

\textbf{(d)} An \emph{isomorphism} between two weighted sets $A$ and $B$ means
a bijection $\rho:A\rightarrow B$ that preserves the weight (i.e., each $a\in
A$ satisfies $\left\vert \rho\left(  a\right)  \right\vert =\left\vert
a\right\vert $). \medskip

\textbf{(e)} We say that two weighted sets $A$ and $B$ are \emph{isomorphic}
(this is written $A\cong B$) if there exists an isomorphism between $A$ and
$B$.
\end{definition}

\begin{proposition}
\label{prop.gf-ws.iso}
\lean{AlgebraicCombinatorics.WeightedSet.weightGenFun_eq_of_isomorphic}
\leantarget
\leanok
Let $A$ and $B$ be two isomorphic finite-type weighted
sets. Then, $\overline{A}=\overline{B}$.
\end{proposition}

\begin{proof}
\leanok
The weighted sets $A$ and $B$ are isomorphic; thus,
there exists an isomorphism $\rho:A\rightarrow B$. Consider this $\rho$. Then,
$\rho$ is a bijection and preserves the weight (since $\rho$ is an isomorphism
of weighted sets). The latter property says that we have
\begin{equation}
\left\vert \rho\left(  a\right)  \right\vert =\left\vert a\right\vert
\ \ \ \ \ \ \ \ \ \ \text{for each }a\in A. \label{pf.prop.gf-ws.iso.wtpres}%
\end{equation}
Now, the definition of $\overline{B}$ yields%
\begin{align*}
\overline{B}  &  =\sum_{b\in B}x^{\left\vert b\right\vert }=\sum_{a\in
A}\underbrace{x^{\left\vert \rho\left(  a\right)  \right\vert }}%
_{\substack{=x^{\left\vert a\right\vert }\\\text{(by the weight-preservation}\\
\text{property above)}}}\ \ \ \ \ \ \ \ \ \ \left(
\begin{array}
[c]{c}%
\text{here, we have substituted }\rho\left(  a\right)  \text{ for }b\\
\text{in the sum, since }\rho\text{ is a bijection}%
\end{array}
\right) \\
&  =\sum_{a\in A}x^{\left\vert a\right\vert }.
\end{align*}
Comparing this with%
\[
\overline{A}=\sum_{a\in A}x^{\left\vert a\right\vert }%
\ \ \ \ \ \ \ \ \ \ \left(  \text{by the definition of }\overline{A}\right)
,
\]
we obtain $\overline{A}=\overline{B}$.
\end{proof}

\begin{definition}
\label{def.gf-ws.djun}
\lean{AlgebraicCombinatorics.WeightedSet.disjointUnion}
\leantarget
\leanok
Let $A$ and $B$ be two weighted sets. Then, the weighted
set $A+B$ is defined to be the disjoint union of $A$ and $B$, with the weight
function inherited from $A$ and $B$ (meaning that each element of $A$ has the
same weight that it had in $A$, and each element of $B$ has the same weight
that it had in $B$). Formally speaking, this means that $A+B$ is the set
$\left(  \left\{  0\right\}  \times A\right)  \cup\left(  \left\{  1\right\}
\times B\right)  $, with the weight function given by
\begin{equation}
\left\vert \left(  0,a\right)  \right\vert =\left\vert a\right\vert
\ \ \ \ \ \ \ \ \ \ \text{for each }a\in A \label{eq.def.gf-ws.djun.wt0}%
\end{equation}
and%
\begin{equation}
\left\vert \left(  1,b\right)  \right\vert =\left\vert b\right\vert
\ \ \ \ \ \ \ \ \ \ \text{for each }b\in B. \label{eq.def.gf-ws.djun.wt1}%
\end{equation}
\end{definition}

\begin{proposition}
\label{prop.gf-ws.djun}
\lean{AlgebraicCombinatorics.WeightedSet.disjointUnion_isFiniteType, AlgebraicCombinatorics.WeightedSet.weightGenFun_disjointUnion}
\leantarget
\leanok
Let $A$ and $B$ be two finite-type weighted sets. Then,
$A+B$ is finite-type, too, and satisfies $\overline{A+B}=\overline
{A}+\overline{B}$.
\end{proposition}

\begin{proof}
\leanok
The formal definition of $A+B$ says that $A+B=\left(  \left\{  0\right\}
\times A\right)  \cup\left(  \left\{  1\right\}  \times B\right)  $. Thus, the
set $A+B$ is the union of the two disjoint sets $\left\{  0\right\}  \times A$
and $\left\{  1\right\}  \times B$ (indeed, these two sets are clearly
disjoint, since $0\neq1$).

Let us first check that $A+B$ is finite-type.

For each $n\in\mathbb{N}$, there are only finitely many $a\in A$ having weight
$\left\vert a\right\vert =n$ (since $A$ is finite-type), and there are only
finitely many $b\in B$ having weight $\left\vert b\right\vert =n$ (since $B$
is finite-type). Hence, there are only finitely many $c\in A+B$ having weight
$\left\vert c\right\vert =n$ (because any such $c$ either has the form
$\left(  0,a\right)  $ for some $a\in A$ having weight $\left\vert
a\right\vert =n$, or has the form $\left(  1,b\right)  $ for some $b\in B$
having weight $\left\vert b\right\vert =n$). In other words, the weighted set
$A+B$ is finite-type.

Let us now check that $\overline{A+B}=\overline{A}+\overline{B}$. Indeed, the
definition of $\overline{A+B}$ yields%
\begin{align*}
\overline{A+B}  &  =\sum_{c\in A+B}x^{\left\vert c\right\vert }\\
&  =\underbrace{\sum_{c\in\left\{  0\right\}  \times A}x^{\left\vert
c\right\vert }}_{\substack{=\sum_{a\in A}x^{\left\vert \left(  0,a\right)
\right\vert }}}+\underbrace{\sum_{c\in\left\{  1\right\}  \times B}x^{\left\vert
c\right\vert }}_{\substack{=\sum_{b\in B}x^{\left\vert \left(  1,b\right)
\right\vert }}}\\
&  =\sum_{a\in A}\underbrace{x^{\left\vert \left(  0,a\right)  \right\vert }%
}_{\substack{=x^{\left\vert a\right\vert }\\\text{(by the definition}\\
\text{of the weight on }A+B\text{)}}}+\sum_{b\in B}\underbrace{x^{\left\vert
\left(  1,b\right)  \right\vert }}_{\substack{=x^{\left\vert b\right\vert
}\\\text{(by the definition}\\
\text{of the weight on }A+B\text{)}}}=\underbrace{\sum_{a\in
A}x^{\left\vert a\right\vert }}_{=\overline{A}}+\underbrace{\sum_{b\in
B}x^{\left\vert b\right\vert }}_{=\overline{B}}=\overline{A}+\overline{B}.
\end{align*}
\end{proof}

\begin{definition}
\label{def.gf-ws.prod}
\lean{AlgebraicCombinatorics.WeightedSet.prod}
\leantarget
\leanok
Let $A$ and $B$ be two weighted sets. Then, the weighted
set $A\times B$ is defined to be the Cartesian product of the sets $A$ and $B$
(that is, the set $\left\{  \left(  a,b\right)  \ \mid\ a\in A\text{ and }b\in
B\right\}  $), with the weight function defined as follows: For any $\left(
a,b\right)  \in A\times B$, we set%
\begin{equation}
\left\vert \left(  a,b\right)  \right\vert =\left\vert a\right\vert
+\left\vert b\right\vert . \label{eq.def.gf-ws.prod.wt}%
\end{equation}
\end{definition}

\begin{proposition}
\label{prop.gf-ws.prod}
\lean{AlgebraicCombinatorics.WeightedSet.prod_isFiniteType, AlgebraicCombinatorics.WeightedSet.weightGenFun_prod}
\leantarget
\leanok
Let $A$ and $B$ be two finite-type weighted sets. Then,
$A\times B$ is finite-type, too, and satisfies $\overline{A\times B}%
=\overline{A}\cdot\overline{B}$.
\end{proposition}

\begin{proof}
\leanok
The proof that
$A\times B$ is finite-type is easy: Let
$n\in\mathbb{N}$. We must prove that there are only finitely many pairs
$\left(  a,b\right)  \in A\times B$ having weight $\left\vert \left(
a,b\right)  \right\vert =n$. Since $\left\vert \left(  a,b\right)  \right\vert
=\left\vert a\right\vert +\left\vert b\right\vert $, these pairs have the
property that $a\in A$ has weight $k$ for some $k\in\left\{  0,1,\ldots
,n\right\}  $, and that $b\in B$ has weight $n-k$ for the same $k$. This
leaves only finitely many options for $k$, only finitely many options for $a$
(since $A$ is finite-type and thus has only finitely many elements of weight
$k$), and only finitely many options for $b$ (since $B$ is finite-type and
thus has only finitely many elements of $n-k$). Altogether, we thus obtain
only finitely many options for $\left(  a,b\right)  $.

Now, the definition of $\overline{A\times B}$ yields%
\begin{align*}
\overline{A\times B}  &  =\sum_{\left(  a,b\right)  \in A\times B}%
\underbrace{x^{\left\vert \left(  a,b\right)  \right\vert }}%
_{\substack{=x^{\left\vert a\right\vert +\left\vert b\right\vert }\\\text{(by the definition}\\
\text{of the weight on }A\times B\text{)}}}\\
&  =\sum_{\left(  a,b\right)  \in A\times B}\underbrace{x^{\left\vert
a\right\vert +\left\vert b\right\vert }}_{=x^{\left\vert a\right\vert }\cdot
x^{\left\vert b\right\vert }}\\
&  =\sum_{\left(  a,b\right)  \in A\times B}x^{\left\vert a\right\vert }\cdot
x^{\left\vert b\right\vert }=\underbrace{\left(  \sum_{a\in A}x^{\left\vert
a\right\vert }\right)  }_{=\overline{A}}\underbrace{\left(  \sum_{b\in
B}x^{\left\vert b\right\vert }\right)  }_{=\overline{B}}=\overline{A}%
\cdot\overline{B}.
\end{align*}
\end{proof}

The weight function on the Cartesian product $A_{1}\times A_{2}\times\cdots\times
A_{k}$ of $k$ weighted sets is defined by%
\begin{equation}
\left\vert \left(  a_{1},a_{2},\ldots,a_{k}\right)  \right\vert =\left\vert
a_{1}\right\vert +\left\vert a_{2}\right\vert +\cdots+\left\vert
a_{k}\right\vert . \label{eq.gf-ws.prodk.weight}%
\end{equation}
As a particular case of such Cartesian products, we obtain Cartesian powers
when we multiply $k$ copies of the same weighted set:

\begin{definition}
Let $A$ be a weighted set. Then, $A^{k}$ (for $k\in\mathbb{N}$) means the
weighted set $\underbrace{A\times A\times\cdots\times A}_{k\text{ times}}$.
\end{definition}

\begin{proposition}
\label{prop.gf-ws.pow}
\lean{AlgebraicCombinatorics.WeightedSet.pow_isFiniteType, AlgebraicCombinatorics.WeightedSet.weightGenFun_pow}
\leantarget
\leanok
Let $A$ be a finite-type weighted set. Let
$k\in\mathbb{N}$. Then, $A^{k}$ is finite-type, too, and satisfies
$\overline{A^{k}}=\overline{A}^{k}$.
\end{proposition}

\begin{proof}
\leanok
By induction on $k$. The base case $k = 0$ is straightforward: $A^0$ consists
of a single element (the empty tuple) with weight $0$, so $\overline{A^0} = 1
= \overline{A}^0$.

For the inductive step, we observe that $A^{k+1} \cong A \times A^k$
as weighted sets (via the isomorphism that splits a $(k{+}1)$-tuple into its
last entry and the remaining $k$-tuple). By Proposition~\ref{prop.gf-ws.iso},
their generating functions are equal. By Proposition~\ref{prop.gf-ws.prod},
$\overline{A \times A^k} = \overline{A} \cdot \overline{A^k}$. By the
induction hypothesis, $\overline{A^k} = \overline{A}^k$. Hence
$\overline{A^{k+1}} = \overline{A} \cdot \overline{A}^k = \overline{A}^{k+1}$.
\end{proof}

\begin{lemma}
\label{lem.gf-ws.pow-succ-equiv}
\lean{AlgebraicCombinatorics.WeightedSet.pow_succ_equiv}
\leanhelper
\leanok
For any weighted set $A$ and any $n \in \mathbb{N}$, there is a canonical
isomorphism of weighted sets $A^{n+1} \cong A \times A^n$.
\end{lemma}

\begin{proof}
\leanok
The isomorphism sends a $(n{+}1)$-tuple $(a_0, a_1, \ldots, a_n)$ to the pair
$(a_n, (a_0, \ldots, a_{n-1}))$. It preserves weight because the sum of entries
is unchanged.
\end{proof}

\begin{lemma}
\label{lem.gf-ws.tuples-finite-type}
\lean{AlgebraicCombinatorics.WeightedSet.tuples_isFiniteType}
\leanhelper
\leanok
Let $A$ be a finite-type weighted set such that every element has weight
$\geq 1$ (i.e., $A$ has positive weights). Then the infinite disjoint union
$A^0 + A^1 + A^2 + \cdots$ (consisting of all finite tuples of elements of
$A$) is also a finite-type weighted set.
\end{lemma}

\begin{proof}
\leanok
If every element has weight $\geq 1$, then a tuple of length $n$ has total
weight $\geq n$. Therefore, for any fixed weight $m$, we only need to consider
tuples of length $\leq m$. For each such length, there are only finitely many
tuples of weight $m$ (since $A^k$ is finite-type by
Proposition~\ref{prop.gf-ws.pow}). Hence the disjoint union is finite-type.
\end{proof}

\subsection{Domino tilings}

\begin{definition}
\label{def.domino.shapes-and-tilings}
\lean{AlgebraicCombinatorics.DominoTilingsZ.Shape, AlgebraicCombinatorics.DominoTilingsZ.Rectangle, AlgebraicCombinatorics.DominoTilingsZ.Domino, AlgebraicCombinatorics.DominoTilingsZ.Tiling, AlgebraicCombinatorics.DominoTilingsZ.numTilings}
\leantarget
\leanok
\textbf{(a)} A \emph{shape} means a
subset of $\mathbb{Z}^{2}$.

We draw each $\left(  i,j\right)  \in\mathbb{Z}^{2}$ as a unit square with
center at the point $\left(  i,j\right)  $ (in Cartesian coordinates); thus, a
shape can be drawn as a cluster of squares. \medskip

\textbf{(b)} For any $n,m\in\mathbb{N}$, the shape $R_{n,m}$ (called the
$n\times m$\emph{-rectangle}) is defined to be%
\[
\left\{  1,2,\ldots,n\right\}  \times\left\{  1,2,\ldots,m\right\}  =\left\{
\left(  i,j\right)  \in\mathbb{Z}^{2}\ \mid\ 1\leq i\leq n\text{ and }1\leq
j\leq m\right\}  .
\]

\textbf{(c)} A \emph{domino} means a size-$2$ shape of the form%
\begin{align*}
&  \left\{  \left(  i,j\right)  ,\ \left(  i+1,j\right)  \right\}  \text{ (a
``\emph{horizontal domino}'')}%
\ \ \ \ \ \ \ \ \ \ \text{or}\\
&  \left\{  \left(  i,j\right)  ,\ \left(  i,j+1\right)  \right\}  \text{ (a
``\emph{vertical domino}'')}%
\end{align*}
for some $\left(  i,j\right)  \in\mathbb{Z}^{2}$. \medskip

\textbf{(d)} A \emph{domino tiling} of a shape $S$ is a set partition of $S$
into dominos (i.e., a set of disjoint dominos whose union is $S$). \medskip

\textbf{(e)} For any $n,m\in\mathbb{N}$, let $d_{n,m}$ be the \# of domino
tilings of $R_{n,m}$.
\end{definition}

We define the weighted set
\[
D := \left\{  \text{domino tilings of height-}2\text{ rectangles}\right\}
= \left\{  \text{domino tilings of }R_{n,2}\text{ with }n\in
\mathbb{N}\right\}  ,
\]
with the weight of a tiling $T$ of $R_{n,2}$ defined to be $\left\vert
T\right\vert := n$. Thus, $D$ is a finite-type weighted set, with weight
generating function $\overline{D}=\sum_{n\in\mathbb{N}}d_{n,2}x^{n}$.

A \emph{fault} of a domino tiling $T$ is a vertical line $\ell$ such that
each domino of $T$ lies either left of $\ell$ or right of $\ell$ (but
does not straddle $\ell$), and
there is at least one domino of $T$ that lies left of $\ell$, and at
least one domino of $T$ that lies right of $\ell$.

A domino tiling is called \emph{faultfree} if it is nonempty and has no
fault. We define the weighted set
\[
F:=\left\{  \text{\textbf{faultfree} domino tilings of }R_{n,2}\text{ with
}n\in\mathbb{N}\right\}
\]
(with the same weights as in $D$).

\begin{lemma}
\label{lem.gf.weighted-set.domino.fd}
\lean{AlgebraicCombinatorics.DominoTilingsZ.tiling_decomposition_isomorphism}
\leantarget
\leanok
Any domino tiling of a height-$2$
rectangle can be decomposed uniquely into a tuple of faultfree tilings of
(usually smaller) height-$2$ rectangles, by cutting it along its faults.
(If the original tiling was faultfree, then it decomposes into a
$1$-tuple. If the original tiling was empty, then it decomposes into a $0$-tuple.)

Moreover, the sum of the weights of the faultfree tilings in the tuple is the
weight of the original tiling. (In other words, if a tiling $T$ decomposes
into the tuple $\left(  T_{1},T_{2},\ldots,T_{k}\right)  $, then $\left\vert
T\right\vert =\left\vert T_{1}\right\vert +\left\vert T_{2}\right\vert
+\cdots+\left\vert T_{k}\right\vert $.)

In the language of weighted sets, this yields an isomorphism
\[
D\cong F^{0}+F^{1}+F^{2}+F^{3}+\cdots,
\]
and therefore, by the infinite analogue of Proposition~\ref{prop.gf-ws.djun},
\[
\overline{D}  =\overline{F^{0}+F^{1}+F^{2}+F^{3}+\cdots}=\overline{F^{0}%
}+\overline{F^{1}}+\overline{F^{2}}+\overline{F^{3}}+\cdots.
\]
By Proposition~\ref{prop.gf-ws.pow}, this equals
\[
\overline{F}^{0}+\overline{F}^{1}+\overline{F}^{2}+\overline{F}^{3}%
+\cdots=\dfrac{1}{1-\overline{F}}.
\]
\end{lemma}

\begin{proof}
The decomposition is defined by cutting the tiling along its faults.
Given a tiling $T$ of $R_{n,2}$:
\begin{itemize}
\item If $n = 0$ (empty tiling), the decomposition is the empty $0$-tuple.
\item If $T$ has no faults, the decomposition is the $1$-tuple $(T)$.
\item If $T$ has faults, let $k$ be the smallest fault position. Cut $T$ at position $k$
to obtain a faultfree left part (a tiling of $R_{k,2}$) and a right part (a tiling of $R_{n-k,2}$).
Recursively decompose the right part, and prepend the left part.
\end{itemize}

The composition (inverse map) concatenates a tuple of faultfree tilings
horizontally by shifting each component to the appropriate position. The
weight-preservation property follows because the total width of the
concatenated rectangle equals the sum of the widths of the components.

The proof that these two maps are inverse to each other requires showing:
\begin{enumerate}
\item Composing and then decomposing recovers the original tuple (by showing
that faults in the composed tiling occur exactly at the partial width sums).
\item Decomposing and then composing recovers the original tiling (by showing
that cutting at faults and reassembling gives back the original domino set).
\end{enumerate}
\end{proof}

\subsubsection{Helpers for the decomposition isomorphism}

\begin{lemma}
\label{lem.gf-ws.domino-left-or-right}
\lean{AlgebraicCombinatorics.DominoTilingsZ.domino_left_or_right_at_fault}
\leanhelper
\leanok
At a fault position $k$ of a tiling $T$ of $R_{n,2}$, each domino of $T$ lies
entirely in the left part (all $x$-coordinates $\le k$) or entirely in the
right part (all $x$-coordinates $\ge k+1$).
\end{lemma}

\begin{proof}
\leanok
For horizontal dominos, the fault condition directly gives that both cells
lie on the same side. For vertical dominos, both cells share the same
$x$-coordinate, so they are trivially on one side.
\end{proof}

\begin{lemma}
\label{lem.gf-ws.restrict-left}
\lean{AlgebraicCombinatorics.DominoTilingsZ.restrictTilingLeft}
\leanhelper
\leanok
Given a tiling $T$ of $R_{n,2}$ with a fault at position $k$, the set of
dominos lying in the left part forms a valid tiling of $R_{k,2}$.
\end{lemma}

\begin{proof}
\leanok
Pairwise disjointness is inherited from $T$. For coverage: every cell
$(x,y)$ with $1 \le x \le k$ is covered by some domino $d \in T$, and
by Lemma~\ref{lem.gf-ws.domino-left-or-right}, $d$ lies entirely in the
left part.
\end{proof}

\begin{lemma}
\label{lem.gf-ws.restrict-right}
\lean{AlgebraicCombinatorics.DominoTilingsZ.restrictTilingRight}
\leanhelper
\leanok
Given a tiling $T$ of $R_{n,2}$ with a fault at position $k$, the set of
dominos lying in the right part, after shifting left by $k$, forms a valid
tiling of $R_{n-k,2}$.
\end{lemma}

\begin{proof}
\leanok
Each right-part domino is shifted by $-k$ in the $x$-coordinate. Disjointness
is preserved since shifting is injective. Coverage follows because every
cell $(x,y)$ with $k+1 \le x \le n$ is covered by a right-part domino.
\end{proof}

\begin{lemma}
\label{lem.gf-ws.restrict-left-faultfree}
\lean{AlgebraicCombinatorics.DominoTilingsZ.restrictTilingLeft_isFaultfree}
\leanhelper
\leanok
The left restriction at the \emph{minimum} fault position is faultfree.
That is, if $k$ is the smallest fault of $T$ and $j < k$, then $j$ is not
a fault of the restricted tiling.
\end{lemma}

\begin{proof}
\leanok
A fault at $j < k$ in the left restriction would induce a fault at $j$ in
the original tiling $T$ (since all horizontal dominos in the left part are
also in $T$). But $k$ is the minimum fault, so no fault at $j < k$ exists
in $T$.
\end{proof}

\begin{lemma}
\label{lem.gf-ws.compose-tilings}
\lean{AlgebraicCombinatorics.DominoTilingsZ.composeTilings}
\leanhelper
\leanok
Given a $k$-tuple of faultfree tilings $(T_1, \ldots, T_k)$ of rectangles
$R_{m_1,2}, \ldots, R_{m_k,2}$, their horizontal concatenation (shifting
each $T_i$ by the cumulative width $m_1 + \cdots + m_{i-1}$) produces a
valid tiling of $R_{m_1 + \cdots + m_k, 2}$.
\end{lemma}

\begin{proof}
\leanok
Dominos from different components have disjoint $x$-coordinate ranges
(Lemma~\ref{lem.gf-ws.compose-component-disjoint}), so they are pairwise
disjoint. Coverage follows because each point in the total rectangle lies
in some component's range and is covered by that component's shifted dominos.
\end{proof}

\begin{lemma}
\label{lem.gf-ws.compose-component-disjoint}
\lean{AlgebraicCombinatorics.DominoTilingsZ.composeTilings_component_dominos_disjoint}
\leanhelper
\leanok
Dominos from different components $i_1 \ne i_2$ in the composition have
disjoint shapes. This follows from the fact that their $x$-coordinate ranges
are separated by the partial width sums.
\end{lemma}

\begin{proof}
\leanok
For $i_1 < i_2$, the partial width sum satisfies
$\sum_{j < i_1} m_j + m_{i_1} \le \sum_{j < i_2} m_j$,
so the $x$-coordinate ranges $[1 + \text{offset}_{i_1}, m_{i_1} + \text{offset}_{i_1}]$
and $[1 + \text{offset}_{i_2}, m_{i_2} + \text{offset}_{i_2}]$ do not overlap.
\end{proof}

\begin{lemma}
\label{lem.gf-ws.decompose-tiling}
\lean{AlgebraicCombinatorics.DominoTilingsZ.decomposeTiling}
\leanhelper
\leanok
Given a tiling $T$ of $R_{n,2}$, there exists a $k$-tuple of faultfree
tilings whose horizontal concatenation reproduces $T$. The decomposition
is constructed recursively: for $n = 0$ it
returns the empty tuple; if $T$ is faultfree it returns the singleton $(T)$;
otherwise it cuts at the minimum fault and recurses on the right part.
\end{lemma}

\begin{proof}
Well-founded recursion on $n$: at each fault $k$, the right part has width
$n - k < n$.
\end{proof}

\begin{lemma}
\label{lem.gf-ws.compose-decompose}
\lean{AlgebraicCombinatorics.DominoTilingsZ.composeTilings_decomposeTiling}
\leanhelper
\leanok
Decomposing a tiling $T$ and then composing the resulting tuple recovers
$T$: $\mathrm{compose}(\mathrm{decompose}(T)) = T$.
\end{lemma}

\begin{proof}
\leanok
By strong induction on $n$. The key step shows that the set of dominos of
$T$ equals the union of the left restriction's dominos and the shifted
right restriction's dominos
(Lemma~\ref{lem.gf-ws.dominos-eq-left-union-shifted-right}),
and by the induction hypothesis the right part reconstructs correctly.
\end{proof}

\begin{lemma}
\label{lem.gf-ws.dominos-eq-left-union-shifted-right}
\lean{AlgebraicCombinatorics.DominoTilingsZ.dominos_eq_left_union_shifted_right}
\leanhelper
\leanok
At a fault position $k$, the original tiling's dominos equal the union of
the left restriction's dominos and the right restriction's dominos shifted
back by $k$.
\end{lemma}

\begin{proof}
\leanok
Every domino is in the left or right part. Left-part dominos are unchanged.
Right-part dominos are shifted by $-k$ and then back by $+k$, recovering
the original.
\end{proof}

\begin{lemma}
\label{lem.gf-ws.decompose-compose}
\lean{AlgebraicCombinatorics.DominoTilingsZ.decomposeTiling_composeTilings}
\leanhelper
\leanok
Composing a $k$-tuple of faultfree tilings and then decomposing recovers
the original tuple: $\mathrm{decompose}(\mathrm{compose}(t_1, \ldots, t_k))
= (t_1, \ldots, t_k)$.
\end{lemma}

\begin{proof}
\leanok
By strong induction on $k$. For $k = 0$ and $k = 1$ the result is immediate.
For $k \ge 2$: the composed tiling has a fault at position $m_1 = |t_1|$
(Lemma~\ref{lem.gf-ws.compose-fault-at-boundary}); no earlier fault exists
(Lemma~\ref{lem.gf-ws.compose-no-fault-before-boundary}); the minimum fault
is $m_1$ (Lemma~\ref{lem.gf-ws.compose-minFault-eq}); the left restriction
equals $t_1$ (Lemma~\ref{lem.gf-ws.restrict-left-compose-eq}); and the right
restriction equals the composition of the tail $(t_2, \ldots, t_k)$
(Lemma~\ref{lem.gf-ws.restrict-right-compose-eq}). The induction hypothesis
on the tail completes the proof.
\end{proof}

\begin{lemma}
\label{lem.gf-ws.compose-fault-at-boundary}
\lean{AlgebraicCombinatorics.DominoTilingsZ.composeTilings_hasFault_at_boundary}
\leanhelper
\leanok
For $k \ge 2$, the composed tiling of a $k$-tuple has a fault at position
$m_1$ (the width of the first component). No horizontal domino crosses this
boundary because dominos from the first component have $x \le m_1$ and
dominos from later components have $x \ge m_1 + 1$.
\end{lemma}

\begin{proof}
\leanok
The width $m_1 > 0$ (since faultfree tilings are nonempty) and $m_1 < m_1 + m_2 + \cdots$
(since $k \ge 2$ and $m_2 > 0$). Each horizontal domino from component $0$
has both $x$-coordinates $\le m_1$; each horizontal domino from component $i \ge 1$
has both $x$-coordinates $\ge m_1 + 1$.
\end{proof}

\begin{lemma}
\label{lem.gf-ws.compose-no-fault-before-boundary}
\lean{AlgebraicCombinatorics.DominoTilingsZ.composeTilings_no_fault_before_boundary}
\leanhelper
\leanok
For $k \ge 1$, if $p < m_1$ (the width of the first component), then $p$ is
not a fault in the composed tiling. This is because a fault at $p$ would
induce a fault at $p$ in the first (faultfree) component.
\end{lemma}

\begin{proof}
\leanok
Every domino from component $0$ that is also in the composed tiling at
position $p$ witnesses the same crossing in the first component. But the
first component is faultfree, so no such fault exists.
\end{proof}

\begin{lemma}
\label{lem.gf-ws.compose-minFault-eq}
\lean{AlgebraicCombinatorics.DominoTilingsZ.composeTilings_minFault_eq}
\leanhelper
\leanok
For $k \ge 2$, the minimum fault position of the composed tiling equals $m_1$
(the width of the first component).
\end{lemma}

\begin{proof}
\leanok
There is a fault at $m_1$ (Lemma~\ref{lem.gf-ws.compose-fault-at-boundary})
and no fault at any $p < m_1$ (Lemma~\ref{lem.gf-ws.compose-no-fault-before-boundary}).
\end{proof}

\begin{lemma}
\label{lem.gf-ws.restrict-left-compose-eq}
\lean{AlgebraicCombinatorics.DominoTilingsZ.restrictTilingLeft_composeTilings_eq}
\leanhelper
\leanok
For $k \ge 2$, the left restriction of the composed tiling at the first
boundary $m_1$ equals the first component's tiling.
\end{lemma}

\begin{proof}
\leanok
Dominos in the left part at position $m_1$ come exactly from component $0$
(since components $i \ge 1$ have $x \ge m_1 + 1$), and the partial width sum
for component $0$ is $0$, so no shifting occurs.
\end{proof}

\begin{lemma}
\label{lem.gf-ws.restrict-right-compose-eq}
\lean{AlgebraicCombinatorics.DominoTilingsZ.restrictTilingRight_composeTilings_eq}
\leanhelper
\leanok
For $k \ge 2$, the right restriction of the composed tiling at the first
boundary $m_1$, after shifting, equals the composition of the tail
$(t_2, \ldots, t_k)$.
\end{lemma}

\begin{proof}
\leanok
The dominos in the right part come from components $i \ge 1$. After
unshifting by $m_1$, each component $i$'s partial width sum decreases by
$m_1$, matching the partial width sums of the tail composition.
\end{proof}

\begin{lemma}
\label{lem.gf-ws.decompose-weight-sum}
\lean{AlgebraicCombinatorics.DominoTilingsZ.decomposeTiling_weight_sum}
\leanhelper
\leanok
The sum of the widths of the faultfree tilings returned by the
decomposition equals the original width $n$. That is,
$\sum_{i} m_i = n$.
\end{lemma}

\begin{proof}
\leanok
Follows from $\mathrm{compose}(\mathrm{decompose}(T)) = T$ (Lemma~\ref{lem.gf-ws.compose-decompose})
and the fact that $\mathrm{compose}$ produces a tiling whose width is the
sum of the component widths.
\end{proof}

The only faultfree tilings of height-$2$ rectangles have width $1$
(one vertical domino) or width $2$ (two horizontal dominos stacked
vertically). Thus, the weighted set $F$ consists of just two tilings: one
of weight $1$ and one of weight $2$. Hence $\overline{F}=x+x^{2}$, and so
\[
\overline{D}=\dfrac{1}{1-\overline{F}}=\dfrac{1}{1-\left(  x+x^{2}\right)
}=\dfrac{1}{1-x-x^{2}}=f_{1}+f_{2}x+f_{3}x^{2}+f_{4}x^{3}+\cdots,
\]
where $\left(  f_{0},f_{1},f_{2},\ldots\right)  $ is the Fibonacci sequence.
Thus $d_{n,2}=f_{n+1}$ for each $n\in\mathbb{N}$.

\begin{lemma}
\label{lem.gf-ws.faultfree-height2-classification}
\lean{AlgebraicCombinatorics.DominoTilingsZ.faultfree_height2_classification}
\leanhelper
\leanok
The only faultfree tilings of height-$2$ rectangles have width $1$ or $2$.
\end{lemma}

\begin{proof}
\leanok
Consider any faultfree tiling of a height-$2$ rectangle $R_{n,2}$ with $n \geq 3$.
Look at the domino that covers the cell $(1,1)$. If it is a vertical domino
$\{(1,1),(1,2)\}$, then this vertical domino constitutes the entire first
column, so there is a fault at position $1$ -- contradiction. If it is a
horizontal domino $\{(1,1),(2,1)\}$, then there must be a horizontal domino
$\{(1,2),(2,2)\}$ stacked above it, and these two dominos cover the first two
columns entirely, so there is a fault at position $2$ -- contradiction.
\end{proof}

\begin{lemma}
\label{lem.gf-ws.weightGenFun-faultfreeHeight2}
\lean{AlgebraicCombinatorics.DominoTilingsZ.weightGenFun_faultfreeHeight2}
\leanhelper
\leanok
The weight generating function of faultfree height-$2$ tilings is $x + x^2$.
\end{lemma}

\begin{proof}
\leanok
By the classification lemma, there is exactly one faultfree tiling of width $1$
(the single vertical domino) and exactly one faultfree tiling of width $2$
(the two horizontal dominos), and no faultfree tilings of any other width.
Hence $\overline{F} = x + x^2$.
\end{proof}

\begin{lemma}
\label{lem.gf-ws.numTilings-recurrence}
\lean{AlgebraicCombinatorics.DominoTilingsZ.numTilings_recurrence}
\leanhelper
\leanok
For any $n \in \mathbb{N}$, $d_{n+2,2} = d_{n,2} + d_{n+1,2}$.
\end{lemma}

\begin{proof}
\leanok
We construct a bijection
$\mathrm{Tiling}(R_{n+2,2}) \simeq \mathrm{Tiling}(R_{n,2}) \sqcup \mathrm{Tiling}(R_{n+1,2})$
by classifying tilings according to what covers the cell $(1,1)$:
\begin{itemize}
\item If a vertical domino $\{(1,1),(1,2)\}$ covers $(1,1)$, then removing it and
shifting gives a tiling of $R_{n+1,2}$.
\item If a horizontal domino $\{(1,1),(2,1)\}$ covers $(1,1)$, then $\{(1,2),(2,2)\}$
must also be present; removing both and shifting gives a tiling of $R_{n,2}$.
\end{itemize}
The inverse maps are given by prepending a vertical domino (resp.\ a pair of
horizontal dominos) and shifting the remaining dominos.
\end{proof}

\subsubsection{Helpers for the Fibonacci recurrence}

\begin{lemma}
\label{lem.gf-ws.exists-domino-covering-11}
For $n \ge 1$, every tiling $T$ of $R_{n,2}$ contains a domino covering
the cell $(1,1)$.
\end{lemma}

\begin{proof}
The cell $(1,1) \in R_{n,2}$ is covered by the tiling. By the coverage
condition, some domino in $T$ contains $(1,1)$.
\end{proof}

\begin{lemma}
\label{lem.gf-ws.domino-covering-11-valid}
If a domino $d$ in a tiling $T$ of $R_{n,2}$ covers $(1,1)$, then $d$ is
either the vertical domino $\{(1,1),(1,2)\}$ or the horizontal domino
$\{(1,1),(2,1)\}$.
\end{lemma}

\begin{proof}
A horizontal domino covering $(1,1)$ must be $\{(1,1),(2,1)\}$ since
extending left to $(0,1)$ would leave the rectangle. A vertical domino
covering $(1,1)$ must be $\{(1,1),(1,2)\}$ since extending down to $(1,0)$
would leave the rectangle.
\end{proof}

\begin{lemma}
\label{lem.gf-ws.horizontal-12-if-horizontal-11}
\lean{AlgebraicCombinatorics.DominoTilingsZ.horizontal_12_if_horizontal_11}
\leanhelper
\leanok
If a tiling $T$ of $R_{n,2}$ (with $n \ge 2$) contains the horizontal domino
$\{(1,1),(2,1)\}$, then $T$ also contains $\{(1,2),(2,2)\}$.
\end{lemma}

\begin{proof}
The cell $(1,2)$ must be covered. The vertical domino $\{(1,1),(1,2)\}$
would overlap with $\{(1,1),(2,1)\}$. So $(1,2)$ must be covered by a
horizontal domino $\{(1,2),(2,2)\}$.
\end{proof}

\begin{lemma}
\label{lem.gf-ws.fault-at-1-if-vertical}
\lean{AlgebraicCombinatorics.DominoTilingsZ.fault_at_1_if_vertical}
\leanhelper
\leanok
If $n \ge 2$ and the tiling $T$ of $R_{n,2}$ contains the vertical domino
$\{(1,1),(1,2)\}$, then $T$ has a fault at position $1$. (The vertical domino
covers the entire first column, so no horizontal domino crosses from column~$1$
to column~$2$.)
\end{lemma}

\begin{proof}
The vertical domino covers both cells of column~$1$. Any other domino
covering a cell in column~$1$ would overlap with the vertical domino.
Hence all other dominos have $x$-coordinates $\ge 2$, giving a fault at $1$.
\end{proof}

\begin{lemma}
\label{lem.gf-ws.fault-at-2-if-horizontal}
\lean{AlgebraicCombinatorics.DominoTilingsZ.fault_at_2_if_horizontal}
\leanhelper
\leanok
If $n \ge 3$ and the tiling $T$ of $R_{n,2}$ contains the horizontal domino
$\{(1,1),(2,1)\}$, then $T$ has a fault at position $2$. The two horizontal
dominos $\{(1,1),(2,1)\}$ and $\{(1,2),(2,2)\}$ cover columns $1$--$2$
entirely.
\end{lemma}

\begin{proof}
By Lemma~\ref{lem.gf-ws.horizontal-12-if-horizontal-11}, $\{(1,2),(2,2)\}$
is also present. These two dominos cover all four cells of columns $1$--$2$.
Any other domino must avoid all these cells, forcing all $x$-coordinates $\ge 3$.
\end{proof}

\begin{lemma}
\label{lem.gf-ws.prepend-vertical}
\lean{AlgebraicCombinatorics.DominoTilingsZ.prependVertical}
\leanhelper
\leanok
Given a tiling $T$ of $R_{m,2}$, prepending a vertical domino $\{(1,1),(1,2)\}$
and shifting all of $T$'s dominos right by $1$ gives a valid tiling of
$R_{m+1,2}$.
\end{lemma}

\begin{proof}
\leanok
Disjointness holds because the vertical domino covers column~$1$ and the
shifted dominos have $x \ge 2$. Coverage: column~$1$ is covered by the
vertical domino; columns $2$ through $m+1$ are covered by the shifted
dominos.
\end{proof}

\begin{lemma}
\label{lem.gf-ws.prepend-horizontal-pair}
\lean{AlgebraicCombinatorics.DominoTilingsZ.prependHorizontalPair}
\leanhelper
\leanok
Given a tiling $T$ of $R_{m,2}$, prepending horizontal dominos
$\{(1,1),(2,1)\}$ and $\{(1,2),(2,2)\}$ and shifting all of $T$'s dominos
right by $2$ gives a valid tiling of $R_{m+2,2}$.
\end{lemma}

\begin{proof}
\leanok
Disjointness holds because the two horizontal dominos cover columns $1$--$2$
(and are disjoint from each other since they cover different rows) and the
shifted dominos have $x \ge 3$. Coverage is verified similarly to
Lemma~\ref{lem.gf-ws.prepend-vertical}.
\end{proof}

\begin{lemma}
\label{lem.gf-ws.tilings-finite}
\lean{AlgebraicCombinatorics.DominoTilingsZ.tilings_finite}
\leanhelper
\leanok
For any $n, m \in \mathbb{N}$, the set of all domino tilings of $R_{n,m}$ is finite.
\end{lemma}

\begin{proof}
\leanok
A tiling is determined by its set of dominos, which is a subset of the
(finite) set of all dominos that fit inside $R_{n,m}$.
\end{proof}

\subsubsection{The Fibonacci formula}

\begin{lemma}
\label{lem.gf-ws.numTilings-2-n}
\lean{AlgebraicCombinatorics.DominoTilingsZ.numTilings_2_n}
\leanhelper
\leanok
For any $n \in \mathbb{N}$, $d_{n,2} = f_{n+1}$, where $(f_k)_{k \ge 0}$ is
the Fibonacci sequence.
\end{lemma}

\begin{proof}
\leanok
By strong induction on $n$. Base cases: $d_{0,2} = 1 = f_1$ and
$d_{1,2} = 1 = f_2$. Inductive step: $d_{n+2,2} = d_{n,2} + d_{n+1,2}$
(Lemma~\ref{lem.gf-ws.numTilings-recurrence}) $= f_{n+1} + f_{n+2}$
(induction hypothesis) $= f_{n+3}$ (Fibonacci recurrence).
\end{proof}

\begin{lemma}
\label{lem.gf-ws.countOfWeight-tilingsHeight2}
\lean{AlgebraicCombinatorics.DominoTilingsZ.countOfWeight_tilingsHeight2_eq}
\leanhelper
\leanok
The count of elements of weight $n$ in the weighted set $D$ of height-$2$
tilings equals $d_{n,2}$.
\end{lemma}

\begin{proof}
\leanok
The elements of weight $n$ in $D$ are exactly the tilings of $R_{n,2}$.
\end{proof}

\begin{lemma}
\label{lem.gf-ws.weightGenFun-tilingsHeight2}
\lean{AlgebraicCombinatorics.DominoTilingsZ.weightGenFun_tilingsHeight2}
\leanhelper
\leanok
The weight generating function of $D$ equals $\sum_n f_{n+1} x^n$, the
Fibonacci generating function.
\end{lemma}

\begin{proof}
\leanok
The $n$-th coefficient is the number of elements of weight $n$ in $D$, which is $d_{n,2} = f_{n+1}$.
\end{proof}
